\chapter{按失效条件分类的反例索引：习题解答}

\begin{solution}{D.1}
周期振荡型取 \(a_n=(-1)^n\)，奇偶子列分别趋于 \(-1,1\)。稠密振荡型可取 \(a_n=\sin n\)；若不引用其稠密性，可改取按区块枚举有理网格的 \([0,1]\)-值数列，使每个有理点出现无限次，它有许多常值收敛子列而全列不收敛。前者由两个固定状态振荡，后者由部分极限众多而失败。
\end{solution}

\begin{solution}{D.2}
取调和部分和 \(H_n=\sum_{k=1}^n1/k\)。相邻差为 \(1/(n+1)\to0\)，但
\[
 H_{2n}-H_n=\sum_{k=n+1}^{2n}\frac1k\ge\frac12.
\]
因此取任意起点 \(N\)，再取 \(n\ge N\)、\(m=2n\)，两项距离至少为 \(1/2\)，不满足 Cauchy 条件。
\end{solution}

\begin{solution}{D.3}
Cauchy 定义只比较值域内部的尾部距离，该有理列完全满足定义。失败的是值域完备性：\(\Q\) 中不存在由这些内部距离确定的极限点。把值域扩大到完备的 \(\R\) 后，同一列收敛到 \(\sqrt2\)。
\end{solution}

\begin{solution}{D.4}
令
\[
 a_{3n}=0,\qquad a_{3n+1}=1,\qquad a_{3n+2}=n.
\]
原列因第三类项无界而无界；指标 \(3n\) 的子列趋于 \(0\)，指标 \(3n+1\) 的子列趋于 \(1\)。
\end{solution}

\begin{solution}{D.5}
无穷远机制：\(f(x)=x^2\) 在 \(\R\) 上，取 \(x_n=n,y_n=n+1/n\)，输入距趋零而像距趋于 \(2\)。缺失边界机制：\(f(x)=1/x\) 在 \((0,1)\) 上，取 \(x_n=1/n,y_n=1/(n+1)\)，输入距趋零而像距恒为 \(1\)。
\end{solution}

\begin{solution}{D.6}
对 \(x,y\ge0\)，
\[
 \abs{\sqrt x-\sqrt y}\le\sqrt{\abs{x-y}},
\]
故取 \(\delta=\varepsilon^2\) 得一致连续。若存在 Lipschitz 常数 \(L\)，令 \(y=0\) 得
\[
 \sqrt x\le Lx,\qquad x>0,
\]
即 \(1/\sqrt x\le L\)，令 \(x\to0+\) 矛盾。
\end{solution}

\begin{solution}{D.7}
定义
\[
 f(x)=
 \begin{cases}
 x,&x<0,\\
 x+1,&x\ge0.
 \end{cases}
\]
它严格递增，在 \(0\) 处左右极限为 \(0,1\)，故不连续。像集为
\[
 (-\infty,0)\cup[1,\infty),
\]
中间区间 \([0,1)\) 未被取得。
\end{solution}

\begin{solution}{D.8}
当 \(x\to+\infty\) 时，\((x+1)-x\to1\)、\(2x-x\to+\infty\)、\(x-2x\to-\infty\)、\((x+\sin x)-x\) 无极限。对 \(0\cdot\infty\)，令无穷因子为 \(x\)，零因子依次取 \(x^{-2},x^{-1},x^{-1/2},\sin x/x\)，乘积依次趋于 \(0,1,+\infty\) 或无极限。
\end{solution}

\begin{solution}{D.9}
\((0,1)\) 有界但缺少端点，覆盖 \((1/n,1)\) 无有限子覆盖；连续函数 \(x\) 不取端点极值。\(\R\) 闭但无界，覆盖 \((-n,n)\) 无有限子覆盖；坏行为来自逃向任意远处。两者分别展示不闭与无界。
\end{solution}

\begin{solution}{D.10}
取有理数列 \(q_n\to\sqrt2\)。它在通常距离下是 Cauchy 列，却不在 \(\Q\) 中收敛，所以 \(\Q\) 不完备。又 \(\sqrt2\) 是 \(\Q\) 在 \(\R\) 中的聚点而不属于 \(\Q\)，故 \(\Q\) 不是 \(\R\) 的闭子集。
\end{solution}

\begin{solution}{D.11}
以 Heine--Cantor 定理为例：删去紧致性，连续函数 \(x^2:\R\to\R\) 不一致连续；失败步骤是无法从无穷多个局部半径中提取有限子覆盖并取正的有限最小值。可恢复结论包括：限制到任意紧集后一致连续；或直接增加全局 Lipschitz/Hölder 估计，也可得到一致连续。
\end{solution}

\begin{solution}{D.12}
符号风险在于后续章节可能采用不同的收敛或拓扑记号；版本风险在于同名定理常有有限测度、局部或全局等不同假设；逻辑风险在于附录若提前引用尚未证明的结论，正文再引用附录，会形成循环。预留位置固定分类框架，同时把正式证明和编号留待依赖章节定稿后补入。
\end{solution}
